\documentclass[11pt]{article}

\usepackage[margin=1in]{geometry}
\usepackage{amsmath,amssymb,bm}
\usepackage{hyperref}

\title{Spherical Harmonic Diagonalization of the Different-Ratio Interface BIE}
\author{KFBIM notes}
\date{\today}

\newcommand{\Gint}{\Omega_{\rm int}}
\newcommand{\Gext}{\Omega_{\rm ext}}
\newcommand{\bi}{\beta_{\rm int}}
\newcommand{\be}{\beta_{\rm ext}}
\newcommand{\bsum}{\beta_{\rm int}+\beta_{\rm ext}}
\newcommand{\lint}{\lambda_{\rm int}}
\newcommand{\lext}{\lambda_{\rm ext}}
\newcommand{\Ylm}{Y_{\ell m}}

\begin{document}
\maketitle

\section*{Setting and conventions}

Let $\Gamma=S^2$ be the unit sphere.  The interior region is the unit ball,
the exterior region is its complement, and the normal $n$ points from
$\Gint$ to $\Gext$.  Jumps use the convention from \texttt{theory.md},
\[
    [u]=\gamma_{\rm int}u-\gamma_{\rm ext}u,
    \qquad
    [\beta\,\partial_n u]
    =
    \bi\,\gamma_{n,{\rm int}}u
    -
    \be\,\gamma_{n,{\rm ext}}u .
\]
After division by $\beta$, each phase solves
\[
    (-\Delta+\lambda_s^2)u_s=q_s,
    \qquad
    \lambda_s^2=\frac{\kappa_s^2}{\beta_s},
    \qquad s\in\{{\rm int},{\rm ext}\}.
\]
Here we specialize to
\[
    \lambda_{\rm int}=0,
    \qquad
    \lambda_{\rm ext}=\lambda>0 .
\]

The formulas below diagonalize the free-space layer potentials on the sphere.
They therefore describe the rotationally invariant model problem.  If the
Green's function is instead the Green's function of a finite rectangular box,
as in a KFBIM FFT solve, the box boundary breaks spherical symmetry and the
operators are not exactly diagonal in spherical harmonics.

\section*{Layer-potential eigenvalues on $S^2$}

Use orthonormal spherical harmonics $\Ylm$, $\ell=0,1,\dots$ and
$-\ell\le m\le\ell$.  For a fixed screening parameter $\lambda_s$, write
$\mathrm{i}_\ell$ and $\mathrm{k}_\ell$ for the modified spherical Bessel
functions of the first and second kind.  Define
\[
    p_\ell^{(s)}
    =
    \lambda_s
    \frac{\mathrm{i}_\ell'(\lambda_s)}
         {\mathrm{i}_\ell(\lambda_s)},
    \qquad
    q_\ell^{(s)}
    =
    \lambda_s
    \frac{\mathrm{k}_\ell'(\lambda_s)}
         {\mathrm{k}_\ell(\lambda_s)} .
\]
For $\lambda_s=0$ these are understood by their Laplace limits,
\[
    p_\ell^{(0)}=\ell,
    \qquad
    q_\ell^{(0)}=-(\ell+1).
\]

For density $\Ylm$, the single-layer potential has boundary value
$\mathsf{s}_\ell^{(s)}\Ylm$, with regular radial factor in the ball and
decaying radial factor outside.  Since
\[
    [\partial_n \mathcal{S}_s\Ylm]
    =
    \partial_n(\mathcal{S}_s\Ylm)_{\rm int}
    -
    \partial_n(\mathcal{S}_s\Ylm)_{\rm ext}
    =
    \Ylm,
\]
one obtains
\[
    \boxed{
    \mathsf{s}_\ell^{(s)}
    =
    \frac{1}{p_\ell^{(s)}-q_\ell^{(s)}} .
    }
\]
Therefore
\[
    S_s\Ylm=\mathsf{s}_\ell^{(s)}\Ylm .
\]

The averaged normal derivative of the single layer gives the adjoint
double-layer operator,
\[
    K_s^*\Ylm=\mathsf{k}_\ell^{(s)}\Ylm,
    \qquad
    \boxed{
    \mathsf{k}_\ell^{(s)}
    =
    \frac{p_\ell^{(s)}+q_\ell^{(s)}}
         {2\bigl(p_\ell^{(s)}-q_\ell^{(s)}\bigr)} .
    }
\]
On the sphere the double-layer operator has the same eigenvalues,
\[
    K_s\Ylm=\mathsf{k}_\ell^{(s)}\Ylm .
\]

For the classical double-layer potential $\mathcal{D}_s$ with
$[\mathcal{D}_s\Ylm]=-\Ylm$, the interior and exterior traces are
\[
    \gamma_{\rm int}\mathcal{D}_s\Ylm
    =
    \left(\mathsf{k}_\ell^{(s)}-\frac12\right)\Ylm,
    \qquad
    \gamma_{\rm ext}\mathcal{D}_s\Ylm
    =
    \left(\mathsf{k}_\ell^{(s)}+\frac12\right)\Ylm .
\]
Its averaged normal derivative is the hypersingular operator
\[
    H_s\Ylm=\mathsf{h}_\ell^{(s)}\Ylm,
    \qquad
    \boxed{
    \mathsf{h}_\ell^{(s)}
    =
    \frac{p_\ell^{(s)}q_\ell^{(s)}}
         {p_\ell^{(s)}-q_\ell^{(s)}} .
    }
\]

For the interior Laplace phase, these reduce to
\[
    \boxed{
    \mathsf{s}_\ell^{({\rm int})}
    =
    \frac{1}{2\ell+1},
    \qquad
    \mathsf{k}_\ell^{({\rm int})}
    =
    -\frac{1}{2(2\ell+1)},
    \qquad
    \mathsf{h}_\ell^{({\rm int})}
    =
    -\frac{\ell(\ell+1)}{2\ell+1}.
    }
\]
For the exterior screened phase,
\[
    \mathsf{s}_\ell^{({\rm ext})}
    =
    \frac{1}{p_\ell^{(\lambda)}-q_\ell^{(\lambda)}},
    \qquad
    \mathsf{k}_\ell^{({\rm ext})}
    =
    \frac{p_\ell^{(\lambda)}+q_\ell^{(\lambda)}}
         {2\bigl(p_\ell^{(\lambda)}-q_\ell^{(\lambda)}\bigr)},
    \qquad
    \mathsf{h}_\ell^{({\rm ext})}
    =
    \frac{p_\ell^{(\lambda)}q_\ell^{(\lambda)}}
         {p_\ell^{(\lambda)}-q_\ell^{(\lambda)}} ,
\]
where $p_\ell^{(\lambda)}$ and $q_\ell^{(\lambda)}$ mean the above
logarithmic derivatives evaluated at $\lambda>0$.

\section*{Different-ratio BIE in each spherical harmonic mode}

Let $a_h$ and $b_h$ denote the forcing-corrected jumps from Section 6.2 of
\texttt{theory.md}.  If there is no volume forcing, then $a_h=a$ and
$b_h=b$.  Expand
\[
    \phi=\sum_{\ell,m}\phi_{\ell m}\Ylm,
    \qquad
    \psi=\sum_{\ell,m}\psi_{\ell m}\Ylm,
\]
\[
    a_h=\sum_{\ell,m}a_{\ell m}\Ylm,
    \qquad
    b_h=\sum_{\ell,m}b_{\ell m}\Ylm .
\]
The Section 6.2 second-kind system is
\[
    \left[
    I-\frac{2\be}{\bsum}K_{\rm int}
      +\frac{2\bi}{\bsum}K_{\rm ext}
    \right]\phi
    +(S_{\rm int}-S_{\rm ext})\psi
    =
    a_h,
\]
\[
    \frac{4\bi\be}{\bsum^2}
    (H_{\rm ext}-H_{\rm int})\phi
    +
    \left[
    I+\frac{2\bi}{\bsum}K_{\rm int}^*
      -\frac{2\be}{\bsum}K_{\rm ext}^*
    \right]\psi
    =
    \frac{2b_h}{\bsum}.
\]
Since every operator is diagonal in $\Ylm$, each mode solves the $2\times 2$
linear system
\[
    \boxed{
    A_\ell \phi_{\ell m}+B_\ell \psi_{\ell m}=a_{\ell m},
    \qquad
    C_\ell \phi_{\ell m}+D_\ell \psi_{\ell m}
    =
    \frac{2b_{\ell m}}{\bsum}.
    }
\]
The coefficients are independent of $m$:
\[
    \boxed{
    A_\ell
    =
    1
    -
    \frac{2\be}{\bsum}\mathsf{k}_\ell^{({\rm int})}
    +
    \frac{2\bi}{\bsum}\mathsf{k}_\ell^{({\rm ext})},
    }
\]
\[
    \boxed{
    B_\ell
    =
    \mathsf{s}_\ell^{({\rm int})}
    -
    \mathsf{s}_\ell^{({\rm ext})},
    }
\]
\[
    \boxed{
    C_\ell
    =
    \frac{4\bi\be}{\bsum^2}
    \left(
    \mathsf{h}_\ell^{({\rm ext})}
    -
    \mathsf{h}_\ell^{({\rm int})}
    \right),
    }
\]
\[
    \boxed{
    D_\ell
    =
    1
    +
    \frac{2\bi}{\bsum}\mathsf{k}_\ell^{({\rm int})}
    -
    \frac{2\be}{\bsum}\mathsf{k}_\ell^{({\rm ext})}.
    }
\]

Equivalently, with
\[
    \Delta_\ell=A_\ell D_\ell-B_\ell C_\ell,
\]
the modal densities are
\[
    \boxed{
    \phi_{\ell m}
    =
    \frac{
    D_\ell a_{\ell m}
    -
    B_\ell\,\frac{2b_{\ell m}}{\bsum}
    }
    {\Delta_\ell},
    \qquad
    \psi_{\ell m}
    =
    \frac{
    -C_\ell a_{\ell m}
    +
    A_\ell\,\frac{2b_{\ell m}}{\bsum}
    }
    {\Delta_\ell}.
    }
\]

\section*{Solvability of every mode}

Assume $\bi>0$, $\be>0$, and $\lambda>0$.  The modal system is solvable if
and only if its homogeneous problem has only the zero solution.  Suppose,
therefore, that $a_{\ell m}=b_{\ell m}=0$ and that
$(\phi_{\ell m},\psi_{\ell m})$ solves the homogeneous modal system.  Form the
corresponding homogeneous layer-potential fields
\[
    u_{\rm int}
    =
    \mathcal{S}_{\rm int}\psi
    -
    \alpha_{\rm int}\mathcal{D}_{\rm int}\phi,
    \qquad
    u_{\rm ext}
    =
    \mathcal{S}_{\rm ext}\psi
    -
    \alpha_{\rm ext}\mathcal{D}_{\rm ext}\phi,
\]
where
\[
    \alpha_{\rm int}=\frac{2\be}{\bsum},
    \qquad
    \alpha_{\rm ext}=\frac{2\bi}{\bsum}.
\]
The two BIE rows are exactly the interface conditions
\[
    \gamma_{\rm int}u_{\rm int}
    =
    \gamma_{\rm ext}u_{\rm ext},
    \qquad
    \bi\,\gamma_{n,{\rm int}}u_{\rm int}
    =
    \be\,\gamma_{n,{\rm ext}}u_{\rm ext}.
\]
The fields satisfy
\[
    -\Delta u_{\rm int}=0
    \quad\hbox{in }B_1,
    \qquad
    (-\Delta+\lambda^2)u_{\rm ext}=0
    \quad\hbox{in }\mathbb{R}^3\setminus\overline{B_1},
\]
and the exterior field decays at infinity.

Multiplying the interior equation by $u_{\rm int}$ and the exterior equation
by $u_{\rm ext}$, integrating by parts, and using the outward normal
conventions gives
\begin{align*}
    &\bi\int_{B_1}|\nabla u_{\rm int}|^2\,dx
    +
    \be\int_{\mathbb{R}^3\setminus B_1}
    \left(|\nabla u_{\rm ext}|^2+\lambda^2|u_{\rm ext}|^2\right)\,dx \\
    &\qquad =
    \int_{S^2}
    \left(
    \bi u_{\rm int}\gamma_{n,{\rm int}}u_{\rm int}
    -
    \be u_{\rm ext}\gamma_{n,{\rm ext}}u_{\rm ext}
    \right)\,dS .
\end{align*}
The boundary integral vanishes because the trace and flux are continuous.
Hence $u_{\rm ext}=0$ and $u_{\rm int}$ is constant.  Continuity with
$u_{\rm ext}=0$ gives $u_{\rm int}=0$.

It remains only to exclude a nonzero null representation of the zero field.
For one spherical harmonic mode, the zero interior trace gives
\[
    \mathsf{s}_\ell^{({\rm int})}\psi_{\ell m}
    -
    \alpha_{\rm int}
    \left(\mathsf{k}_\ell^{({\rm int})}-\frac12\right)
    \phi_{\ell m}
    =
    0,
\]
or, using the Laplace eigenvalues,
\[
    \psi_{\ell m}
    =
    -\alpha_{\rm int}(\ell+1)\phi_{\ell m}.
\]
The zero exterior trace gives
\[
    \mathsf{s}_\ell^{({\rm ext})}\psi_{\ell m}
    -
    \alpha_{\rm ext}
    \left(\mathsf{k}_\ell^{({\rm ext})}+\frac12\right)
    \phi_{\ell m}
    =
    0.
\]
Since
\[
    \mathsf{k}_\ell^{({\rm ext})}+\frac12
    =
    \frac{p_\ell^{(\lambda)}}
         {p_\ell^{(\lambda)}-q_\ell^{(\lambda)}}
    =
    p_\ell^{(\lambda)}\mathsf{s}_\ell^{({\rm ext})},
\]
this becomes
\[
    \psi_{\ell m}
    =
    \alpha_{\rm ext}p_\ell^{(\lambda)}\phi_{\ell m}.
\]
Here $p_\ell^{(\lambda)}>0$, while
$\alpha_{\rm int},\alpha_{\rm ext}>0$ and $\ell+1>0$.  The two displayed
relations are compatible only when $\phi_{\ell m}=0$, and then
$\psi_{\ell m}=0$.  Thus the homogeneous modal system has no nonzero solution,
so
\[
    \boxed{\Delta_\ell\ne 0\qquad\hbox{for every }\ell=0,1,2,\dots .}
\]

\section*{Large-mode asymptotics}

Let
\[
    N_\ell=2\ell+1.
\]
The asymptotics below are for fixed exterior screening parameter
$\lambda>0$ as $\ell\to\infty$.  The regular and decaying logarithmic
derivatives satisfy
\[
    p_\ell^{(\lambda)}
    =
    \ell+\frac{\lambda^2}{N_\ell+2}+O(N_\ell^{-3}),
    \qquad
    q_\ell^{(\lambda)}
    =
    -(\ell+1)-\frac{\lambda^2}{N_\ell-2}
    +O(N_\ell^{-3}).
\]
Consequently,
\[
    \mathsf{s}_\ell^{({\rm ext})}
    =
    \frac{1}{N_\ell}
    -
    \frac{2\lambda^2}{N_\ell(N_\ell^2-4)}
    +O(N_\ell^{-5}),
\]
\[
    \mathsf{k}_\ell^{({\rm ext})}
    =
    -\frac{1}{2N_\ell}
    -
    \frac{\lambda^2}{N_\ell(N_\ell^2-4)}
    +O(N_\ell^{-4}),
\]
and
\[
    \mathsf{h}_\ell^{({\rm ext})}
    -
    \mathsf{h}_\ell^{({\rm int})}
    =
    -\frac{\lambda^2(N_\ell^2-3)}
           {2N_\ell(N_\ell^2-4)}
    +O(N_\ell^{-3}).
\]

Substituting these into the modal matrix gives
\[
    \boxed{
    A_\ell
    =
    1
    +
    \frac{\be-\bi}{\bsum\,N_\ell}
    -
    \frac{2\bi\lambda^2}
         {\bsum\,N_\ell(N_\ell^2-4)}
    +O(N_\ell^{-4}),
    }
\]
\[
    \boxed{
    B_\ell
    =
    \frac{2\lambda^2}
         {N_\ell(N_\ell^2-4)}
    +O(N_\ell^{-5}),
    }
\]
\[
    \boxed{
    C_\ell
    =
    -
    \frac{2\bi\be\lambda^2}{\bsum^2}
    \frac{N_\ell^2-3}{N_\ell(N_\ell^2-4)}
    +O(N_\ell^{-3}),
    }
\]
\[
    \boxed{
    D_\ell
    =
    1
    +
    \frac{\be-\bi}{\bsum\,N_\ell}
    +
    \frac{2\be\lambda^2}
         {\bsum\,N_\ell(N_\ell^2-4)}
    +O(N_\ell^{-4}).
    }
\]
In particular,
\[
    B_\ell=O(\ell^{-3}),
    \qquad
    C_\ell=O(\ell^{-1}),
    \qquad
    A_\ell,D_\ell=1+O(\ell^{-1}),
\]
and the determinant obeys
\[
    \Delta_\ell
    =
    1
    +
    \frac{2(\be-\bi)}{\bsum\,N_\ell}
    +O(N_\ell^{-2}).
\]
Thus the modal matrices become identity plus a compact lower-order
perturbation at high spherical harmonic degree.

\section*{Remarks}

The leading hypersingular pieces cancel in the same way as in the general
smooth-interface derivation.  On the sphere this cancellation is visible mode
by mode: $H_{\rm ext}-H_{\rm int}$ is represented only by the difference of
the scalar eigenvalues $\mathsf{h}_\ell^{({\rm ext})}$ and
$\mathsf{h}_\ell^{({\rm int})}$, while the prefactor
$4\bi\be/\bsum^2$ is bounded.

The final implementation paragraph in Section 6.2 of \texttt{theory.md} uses
the KFBIM double-layer jump primitive $[u]=\phi$, which is the opposite of the
classical double-layer sign used above.  The modal diagonalization is
unchanged; only the signs in the double-layer rows should be replaced by the
implementation variant written there.

\end{document}
